% Chapter 2 revised: critical literature synthesis, August 2026
\label{ch:theory}

\noindent
The central problem addressed in this chapter is not whether health is multidimensional. That proposition is well established. Nor is it whether health can be represented by a composite score. The literature contains several successful examples of multidimensional measurement, including the Global Burden of Disease framework, health-expectancy measures, frailty indices, intrinsic-capacity measures, capability measures and latent health metrics developed through factor analysis and item-response theory. The more difficult question is what is gained, and what is lost, when these traditions are brought together in a single construct intended to represent an individual's health-related potential rather than a current health state, a count of deficits, or a probability of an adverse event.

The chapter therefore treats the literature as an argument rather than a catalogue. Five bodies of scholarship are placed in dialogue: (i) health-state and health-loss measurement; (ii) frailty and deficit accumulation; (iii) the WHO healthy-ageing and intrinsic-capacity framework; (iv) capability theory and capability measurement; and (v) social determinants and social vulnerability. The purpose of bringing these traditions together is not to imply that they are interchangeable. On the contrary, their differences provide the theoretical problem that the thesis must resolve. If intrinsic capacity describes what an individual is able to draw upon, functional ability describes what is realised in context, capability describes the opportunity to achieve valued states and activities, and social vulnerability describes accumulated exposure to social disadvantage, then a proposed health-capability construct must specify precisely where it sits among these concepts.

This leads to a deliberately restrictive claim. The $H_{\text{potential}}$ is not proposed as the first multidimensional health measure, the first latent health score, or the first measure to incorporate social conditions. ATHLOS, for example, constructed a longitudinal health metric using exploratory and confirmatory factor analysis followed by Bayesian multilevel item-response theory, while the WHO framework explicitly links intrinsic capacity, environment and functional ability. \cite{caballero_athlos_2017,who_2015_ageing,who_healthy_ageing_2026} Capability measures such as ICECAP-O likewise demonstrate that capability can be operationalised at the individual level, although the measure is intentionally broader than health. \cite{icecap_o,icecap_o_validation} The proposed contribution is narrower: to test whether a health-bounded construct can explicitly connect biological and functional capacity with enabling social conditions and structural constraints, and whether that construct remains empirically distinguishable from established measures of intrinsic capacity, frailty, social vulnerability and current health.

The chapter consequently ends not with a declaration that the proposed construct is valid, but with a set of propositions that the empirical work must attempt to falsify. The logic is important. A thesis of this kind is strongest when the data are allowed to demonstrate that the proposed construct is unnecessary, multidimensional without a valid general factor, or prognostically redundant. The objective is not to protect the proposed index from such findings but to determine what the evidence can legitimately support.

\section{What Does It Mean to Measure Health?}

Measurement begins with a construct, not an algorithm. A score can be calculated whenever observations are combined, but calculation does not establish that the resulting quantity corresponds to a meaningful latent attribute. This distinction is particularly important in multidimensional health research because the word \emph{health} is used to describe states, capacities, functionings, opportunities, vulnerabilities and outcomes that occupy different conceptual levels.

The WHO healthy-ageing framework provides an important starting point because it explicitly distinguishes intrinsic capacity from functional ability. Intrinsic capacity refers to the composite of an individual's physical and mental capacities, while functional ability emerges from the interaction between those capacities and relevant environmental characteristics. The environment is defined broadly, including physical, social and policy contexts and the systems and relationships through which people live. \cite{who_2015_ageing,who_healthy_ageing_2026} This formulation prevents a common measurement error: interpreting an observed limitation as evidence that the underlying biological capacity is absent when the observed limitation may instead reflect an environmental barrier.

The distinction can be expressed conceptually, for person $i$, as
\begin{equation}
\boxed{FA_i=f(IC_i,E_i,IC_i\times E_i)},
\label{eq:who_fa_conceptual}
\end{equation}
where $FA_i$ denotes realised functional ability, $IC_i$ intrinsic capacity and
$E_i$ the relevant environment. Equation~\ref{eq:who_fa_conceptual} is a
translation of the WHO conceptual architecture into notation; it is not an
equation estimated or endorsed by WHO. Writing the interaction among the
arguments of $f$ does not assert a linear interaction term, a multiplicative
production function or any particular scale.

The conceptual sequence used in this thesis is
\begin{equation}
\boxed{\begin{aligned}
\text{health conditions}&\longrightarrow\text{intrinsic capacities}
\longrightarrow\text{health-related capabilities}\\
&\longrightarrow\text{functional ability}
\longrightarrow\text{observed outcomes}
\end{aligned}},
\label{eq:conceptual_sequence}
\end{equation}
with environmental and social structures acting across the sequence. The
arrows express theoretical ordering, not identified causal effects.

The thesis therefore asks whether underlying health-related potential can be
measured separately from realised function. Its general construct statement is
\begin{equation}
\boxed{H_{P,i}=f(BP_i,SM_i,SB_i)},
\label{eq:hp_general_class}
\end{equation}
where $BP_i$ is biological and functional capacity, $SM_i$ comprises social and
environmental enabling conditions, and $SB_i$ comprises structural and social
barriers. The mapping $IC\rightarrow BP$ is approximate rather than an asserted
equivalence. Likewise, $E\rightarrow(SM,SB)$ separates enabling conditions from
constraints without implying that every environmental variable belongs inside
the construct. Whether social structure is a component, determinant or modifier
is tested through competing specifications.

\begin{table}[htbp]
\centering
\caption{Conceptual and empirical roles in the H-Potential argument}
\label{tab:concept_roles}
\small
\begin{tabular}{p{0.22\linewidth}p{0.22\linewidth}p{0.44\linewidth}}
\hline
Concept & Mathematical representation & Empirical role \\
\hline
Intrinsic capacity & $IC$ & Latent or domain-specific capacities; principal clinical predecessor \\
Environment & $E$ & Context, enablement or constraint; not presumed reflective \\
Functional ability & $FA=f(IC,E,IC\times E)$ & Realised capability and functioning \\
Frailty & $F$ & Competing formative/composite construct \\
Social vulnerability & $SV$ & Contextual competitor and possible determinant \\
H-Potential & $H_P=f(BP,SM,SB)$ & Proposed health-bounded construct \\
Mortality & $(T,D)$ & Downstream prognostic outcome, never a defining indicator \\
Hospitalisation & $(T_H,D_H)$ & Downstream outcome with death as a competing event \\
QALY & utility $\times$ survival & Cardinal economic comparator, not a latent-score interpretation \\
Health expectancy & survival/function & Population-level comparator \\
\hline
\end{tabular}
\end{table}

The capability approach provides a second conceptual bridge. Sen distinguishes achieved functionings from the capability to achieve alternative functionings. \cite{sen_1993} This distinction is useful in health because the same observed outcome can arise from different opportunity structures. A person may not participate in an activity because of choice, because of impaired capacity, or because environmental and institutional constraints make participation difficult. A measurement system that observes only achieved functioning can therefore conflate inability with constrained opportunity.

However, the capability approach cannot be imported into health measurement without qualification. Capability is a theory of human freedom and wellbeing, not a ready-made psychometric definition of health. ICECAP-O illustrates this boundary explicitly: it was developed as a measure of capability wellbeing in older people and intentionally extends beyond health states. \cite{icecap_o,icecap_o_validation} Empirical comparisons also suggest that ICECAP-O and direct health measures overlap only partially, particularly in physical functioning. \cite{icecap_o_validation} This is not a weakness of ICECAP-O; it is evidence that capability and health are related constructs with different scopes.

The present thesis therefore uses the term \emph{health-related capability} in a deliberately restricted sense. It refers to the capacity and opportunity to sustain health-related functioning under the biological and social conditions in which an individual lives. It is not intended to measure total human capability, flourishing, life satisfaction or general wellbeing. This boundary is necessary for discriminant validity: if the construct can only be distinguished from general wellbeing by its name, its theoretical specification is inadequate.

A further distinction is required between a latent construct and a prognostic outcome. A score derived from NHATS is not itself future health. It is an estimate of a latent state hypothesised to be associated with subsequent outcomes. Mortality and hospitalisation therefore cannot be used to define the construct without creating circularity. The measurement model must be developed before the prognostic model is evaluated. Prediction is evidence relevant to validity, but prediction alone does not establish construct meaning.

\section{From Health Loss to Health Potential: QALYs, DALYs and Health Expectancy}

QALYs and DALYs established the importance of constructing common measures across heterogeneous health states. QALYs combine duration and health-related quality, while DALYs combine years of life lost and years lived with disability. Their principal achievement is comparability: health consequences that arise from different diseases can be represented on a common evaluative scale. \cite{whitehead_2010,drummond2015,sundmacher_2017}

Their conceptual purpose, however, differs from that of the present thesis. QALYs and DALYs are primarily measures of health gain, health loss and population burden. Their construction requires explicit valuation choices concerning health states, disability and, in some applications, time. These choices are not neutral technical operations. They encode assumptions about what counts as a loss and how different states should be compared. \cite{drummond2015}

The economic case for these measures should not be understated. A common unit
permits opportunity costs to be discussed across diseases and makes otherwise
implicit allocation choices open to scrutiny. The objection is not that QALYs
are too simple and should be displaced by a richer score. It is that their
question is different. A QALY is an evaluative consequence measure; it does not
claim to identify the biological reserve, social opportunity and institutional
conditions from which a health state arose. Nor does an H-Potential factor score
have the preference-based or ratio-scale properties needed to replace a QALY.

Capability-based criticism adds a second qualification. Cookson showed that the
capability approach can inform health evaluation in several ways, but direct
estimation and valuation of capability sets presents serious empirical
difficulties; standard health-state instruments may remain useful when the
health consequences of an intervention are the decision-maker's principal
concern. \cite{cookson_2005} H-Potential therefore complements rather than
supersedes economic evaluation.

Health expectancy moves closer to the present problem because it introduces time directly. Healthy life expectancy and related measures combine mortality and morbidity or disability information to estimate expected years lived in a specified health state. \cite{who_hale_2019} Yet health expectancy is principally a population-level demographic quantity. It is particularly valuable for describing inequalities between populations and changes over time, but it does not necessarily provide an individual-level latent representation of the mechanisms that generate future health.

This creates an important conceptual bridge. Health expectancy asks, in effect, \emph{how many years of life are expected to be lived in a defined health state?} A prognostic risk model asks \emph{what is the probability of a specified event?} The proposed construct asks a different question: \emph{what level of health-related capacity and opportunity is available now, given the person's biological and social context?} The empirical relationship between that state and future outcomes must then be demonstrated rather than built into the definition.

The distinction prevents a recurring problem in composite-index development: a predictive score can perform well while remaining theoretically uninterpretable. Conversely, a theoretically coherent construct may predict only modestly because its purpose is broader than any single outcome. The thesis therefore treats prognostic performance as one form of validity evidence rather than the criterion by which the construct is created.

\section{Quality of Life, Functioning and the Problem of the Observed State}

Health-related quality-of-life instruments provide an indispensable account of experienced health. Generic measures such as EQ-5D and disease-specific measures such as Haem-A-QoL capture symptoms, functioning and perceived impact with a degree of clinical specificity that a broad latent construct cannot reproduce. Their value is precisely that they remain close to the lived health state.

The limitation for the present purpose is conceptual rather than methodological. A current health-state instrument does not necessarily distinguish between the capacity to function and the circumstances that permit that capacity to be realised. This is especially relevant to chronic disease, where treatment can preserve biological function while social circumstances determine whether that function can be translated into participation, access, independence and continuity of care.

The haemophilia literature makes this distinction particularly important. Modern management has substantially altered the natural history of haemophilia, making survival increasingly compatible with long-term management of musculoskeletal damage, treatment burden, mental health, comorbidity and access to specialist care. \cite{mauser_2009,srivastava_2020} A disease-specific quality-of-life measure can capture consequences of these processes, but it should not be expected to serve simultaneously as a model of biological reserve, social opportunity and structural constraint.

This gives rise to a testable requirement for $H_{\text{potential}}$: it should correlate with established health-state measures because they share substantive information, but it should not be reducible to them. If the proposed score explains no information beyond current function, its claim to represent \emph{potential} is weakened. If it is almost independent of established health measures, its health content is also questionable. Convergent and discriminant validity therefore need to be considered together.

\section{Frailty: A Precedent for Multidimensional Vulnerability}

Frailty provides perhaps the strongest precedent for the thesis because it demonstrates that heterogeneous observations can be aggregated into a meaningful state variable with prognostic consequences. The deficit-accumulation model treats symptoms, diseases, disabilities and other deficits as accumulated manifestations of vulnerability. The resulting frailty index has a characteristic distribution and predicts adverse outcomes. \cite{rockwood_mitnitski_2007}

The intellectual importance of deficit accumulation is not merely that it produces a score. It demonstrates that a high-dimensional representation of health can possess regular empirical properties even when individual deficits are heterogeneous. Earlier work also showed that different sets of deficits can produce broadly comparable prognostic information, suggesting that the construct captures a general state of vulnerability rather than the effect of a single disease. \cite{rockwood_mitnitski_2007}

Yet frailty does not solve the problem posed here. Frailty is principally concerned with vulnerability and reduced reserve. The proposed construct is intended to represent both capacity and the conditions under which capacity can be exercised. Moreover, deficit accumulation generally treats observed deficits as components of vulnerability. Social disadvantage may appear as a deficit, but the framework does not necessarily distinguish between a biological limitation and an external barrier that constrains its expression.

A fair counterargument is that this distinction may not justify another score.
Frailty already summarises vulnerability efficiently, predicts clinically
important outcomes and is familiar to geriatric practice. A parsimonious
frailty model may be preferable to a theoretically richer index that is harder
to calculate, explain or act upon. Deficit accumulation is deliberately
agnostic about the causal origin of each deficit; that apparent theoretical
limitation contributes to its robustness across different deficit sets.

H-Potential can answer this objection only through evidence. It must show either
that two people with similar frailty can have meaningfully different
health-related opportunity because their enabling conditions differ, or that
the domain profile changes interpretation, prediction or intervention planning.
If it merely reverses the sign of frailty and appends socioeconomic variables,
parsimony favours frailty.

This distinction creates an empirical challenge rather than a rhetorical one. If the H-potential score is simply a positively oriented frailty index, the conceptual contribution is limited. The analysis must therefore compare the proposed score against a pre-specified frailty measure and determine whether it captures information not already contained in deficit accumulation. This comparison should be performed both at the measurement level and at the outcome level.

The comparison also exposes a deeper measurement issue. A frailty index is commonly formative or composite in logic: the deficits contribute to the state rather than being interchangeable manifestations of a single latent trait. A conventional reflective CFA, in which every observed indicator is assumed to be an effect of one common factor, is therefore not automatically appropriate for all health constructs. The thesis must not choose a measurement model because SEM makes it convenient; it must choose the model that best represents the ontology of the construct.

\section{Intrinsic Capacity: From Conceptual Framework to Measurement Problem}

WHO's intrinsic-capacity framework provides the most direct biological predecessor to the proposed Biomedical Potential domain. Intrinsic capacity encompasses multiple physical and mental capacities, conventionally organised into locomotor, cognitive, psychological, sensory and vitality domains. \cite{who_intrinsic_capacity_2017,beard_intrinsic_2019}

The framework is theoretically powerful because it shifts attention from diagnosed disease to what a person is able to draw upon. It also provides a direct bridge to functional ability: declining intrinsic capacity need not translate into equivalent loss of functioning when environments are supportive. \cite{who_2015_ageing,who_healthy_ageing_2026} This is precisely the conceptual space in which a health-capability construct might reside.

However, the intrinsic-capacity literature contains an important unresolved measurement problem. A scoping review of 31 studies found substantial inconsistency in the conceptualisation and operationalisation of intrinsic capacity, with studies implicitly or explicitly treating the construct as both reflective and formative. \cite{who_intrinsic_capacity_2017,beard_intrinsic_2019} This is not a minor technical disagreement. If intrinsic capacity is reflective, observed capacities are treated as manifestations of an underlying common trait. If it is formative, the capacities jointly define the construct and may change independently without implying that the latent state has changed in the same direction.

The empirical literature also demonstrates heterogeneity in how the five domains are measured. A later scoping review identified 53 studies and found considerable variation in both domain selection and score construction, with sensory and vitality domains particularly heterogeneous. \cite{intrinsic_capacity_predictive_2023} This heterogeneity creates both a limitation and an opportunity for the present research. The limitation is that a proposed H-potential score cannot assume that a single canonical operationalisation of intrinsic capacity exists. The opportunity is that the thesis can make the measurement assumptions explicit and test competing structures.

Longitudinal evidence provides an additional bridge. Intrinsic-capacity measures have been associated with subsequent disability, hospital admission and mortality. \cite{beard_intrinsic_2019,intrinsic_capacity_predictive_2023} Such findings support the prognostic relevance of capacity but do not establish that intrinsic capacity is equivalent to a broader health-capability construct. Indeed, if intrinsic capacity already predicts the outcomes of interest strongly, the burden of proof shifts to H-potential: it must demonstrate incremental information rather than merely reproduce the same association.

The most direct criticism comes from within gerontology. Hoogendijk and
colleagues argue that revisions to the domain structure, heterogeneous
operationalisations, uncertain reflective-versus-formative logic and limited
independent validation have made intrinsic capacity difficult to interpret.
They further warn that studies often measure the absence of familiar deficits
rather than capacity itself, risking the relabelling of established knowledge.
\cite{hoogendijk_2023} Advocates respond that a capacity-centred framework
corrects limitations of disease-counting and can organise prevention around
functions that matter in later life, while acknowledging that interventions to
optimise intrinsic capacity still require testing. \cite{geerinck_2022}

This debate makes intrinsic capacity both the closest clinical framework and a
cautionary case. H-Potential closes no gap merely by adding context to an
uncertain total score. Its proposed advance is procedural: define the domain
ontology before estimation; separate capacity indicators from conditions that
convert capacity into function; compare reflective, formative and profile
interpretations; and permit the analysis to reject a general score.

\section{Multidimensional Health Metrics: ATHLOS and the Problem of Novelty}

The ATHLOS project is a particularly important precedent because it demonstrates that a single longitudinal health metric can be constructed from heterogeneous indicators using modern latent-variable methods. Caballero et al. used exploratory and second-order confirmatory factor analysis to identify a multidimensional structure and then applied Bayesian multilevel item-response theory to create a health score that could be compared across waves. \cite{caballero_athlos_2017}

This work directly challenges any claim that the proposed thesis is novel merely because it uses factor analysis, IRT or longitudinal health measurement. ATHLOS also demonstrates that the move from many indicators to one score is not trivial: the investigators first investigated dimensionality and then calibrated a score using item-response methods. The methodological lesson is important for H-potential. A general score must be earned by evidence of common structure, not assumed because a scalar index is convenient.

At the same time, ATHLOS and H-potential address different conceptual questions. ATHLOS operationalised health as a multidimensional state of functioning and used a large set of indicators to derive a general health metric. The present thesis is more explicitly concerned with the relationship between biological capacity, enabling social conditions and structural constraints. The distinction must therefore be tested rather than asserted. If the proposed domains collapse empirically into a general health factor that is indistinguishable from established health metrics, the novelty claim should be narrowed accordingly.

The appropriate lesson from ATHLOS is therefore methodological humility. The thesis should adopt the successful part of the precedent---explicit investigation of dimensionality and longitudinal comparability---while resisting the temptation to treat any general factor as evidence of a new construct. Construct identity depends on the content and causal interpretation of the indicators, not on the existence of a statistically adequate one-factor model.

\section{Capability Theory and the Boundary Between Health and Wellbeing}

Capability theory supplies the normative bridge between capacity and realised functioning. Sen's central distinction is between what people actually achieve and the substantive opportunities available to them. \cite{sen_1993} This is particularly relevant when health is measured in social context because observed functioning can be constrained by factors external to the individual.

ICECAP-O provides the clearest operational precedent. It was developed through qualitative work with older people and was designed to capture broader capability wellbeing rather than health alone. \cite{icecap_o} Subsequent validation has shown associations with both health and wellbeing, while comparisons with direct health measures indicate incomplete overlap. \cite{icecap_o_validation} The evidence therefore supports the conceptual usefulness of capability measurement while simultaneously warning against collapsing capability wellbeing into health status.

For H-potential, the implication is that capability theory should function as a bridge, not as a replacement for health measurement. The proposed construct must remain health-bounded. Its social domains should represent conditions relevant to the maintenance or realisation of health-related functioning, rather than every determinant of a good life.

There is, however, a fundamental empirical difficulty. Functionings can often
be observed, but capability is a set of genuine opportunities, including options
not chosen. The same non-participation may represent preference, temporary
circumstance, coercion or inability. Inferring an opportunity set from observed
behaviour therefore risks calling constrained adaptation a free choice, while
self-report can be affected by expectations and adaptation. \cite{cookson_2005}

H-Potential does not claim to observe the complete capability set. It makes the
more limited claim that selected biological capacities and documented enabling
or constraining conditions contain information about health-related opportunity.
The score is an estimate under a stated measurement model, not a direct reading
of human freedom. Qualitative work and content validation remain necessary
because covariance cannot establish that the selected opportunities are those
people value.

This distinction also has consequences for indicator selection. Variables such as social participation or perceived security may be relevant because they influence health-related capability, but their inclusion requires a theoretical pathway. They should not be included simply because they correlate with wellbeing. The code should therefore maintain a transparent mapping from each indicator to the theoretical proposition it is intended to operationalise.

\section{Social Determinants and the Conversion of Capacity into Function}

The social-determinants literature supplies the contextual bridge that biological models alone cannot provide. The WHO framework distinguishes structural determinants and socioeconomic position from intermediary determinants, emphasising that individual health is patterned by the social organisation of resources and opportunities. \cite{solar_irwin_2010}

This matters because a health-capability construct can otherwise commit an attribution error. A person may possess adequate physical capacity but be unable to use it because of transport barriers, unsafe housing, inadequate financial resources, limited access to care or social isolation. These are not simply additional biomedical deficits. They are conditions that alter the conversion of capacity into functioning.

The life-course perspective strengthens this argument. Social advantage and disadvantage accumulate through education, employment, income, housing, exposure and access to resources. Consequently, an observed social condition in later life can be both a current constraint and the downstream expression of earlier exposures. The relevance to H-potential is therefore temporal: the construct is measured at a point in time but should be interpreted as the current expression of a process extending across the life course.

Empirical evidence supports the importance of these pathways. Social determinants have been associated with mortality among older adults, and ATHLOS analyses have demonstrated relationships among social determinants, health status and long-term mortality. \cite{social_determinants_mortality_2014,athlos_social_determinants_2018} These findings justify the inclusion of social conditions as theoretically relevant information. They do not, however, establish that social variables should be reflective indicators of a latent health factor. Their causal or contextual status must remain explicit.

Including social conditions inside an individual score creates a serious
counterargument. It may conflate the person with the conditions imposed upon
them. A low score could then be read as an individual deficiency when the
appropriate conclusion is that housing, transport, income or service access is
inadequate. It can also hide policy-relevant information: two people can receive
the same total through very different combinations of biological limitation and
structural constraint.

The theory addresses this by preserving domain scores and treating the scalar
score as conditional rather than essential. Biomedical Potential describes
capacity-related information; Social Modifiers describe enablement; Systemic
Burdens describe constraints. A total is reported only if common structure and
an interpretable aggregation rule are supported. Even then, the profile remains
available so that structural burden is never silently reclassified as personal
pathology.

\section{Social Vulnerability: The Most Direct Competitor}

Social vulnerability is perhaps the most important comparator because it already aggregates adverse social circumstances into an individual-level construct and has demonstrated prognostic relevance. Andrew, Mitnitski and Rockwood showed that social vulnerability was associated with mortality in older adults, while subsequent SHARE analyses demonstrated associations with mortality and disability across countries. \cite{andrew_social_vulnerability_2008,social_vulnerability_share_2014}

The existence of this literature changes the burden of proof for H-potential. It is no longer sufficient to argue that social conditions matter. They clearly do. The question is whether the proposed Social Modifiers and Systemic Burdens represent something more specific than social vulnerability, and whether their integration with biological capacity adds information that neither component provides independently.

This suggests a three-way empirical comparison. First, does H-potential differ from biomedical capacity measures? Second, does it differ from social vulnerability? Third, does it retain information when both are considered? If the answer to the third question is no, the proposed integrated construct may have little justification as a separate metric even if its individual domains are meaningful.

The comparison should also consider the direction of interpretation. Social vulnerability is generally framed as accumulated disadvantage and risk. H-potential is intended to be framed positively as available health-related potential, while systemic burdens reduce that potential. These are not merely opposite labels. A positive transformation of a risk score is not automatically a new construct. The thesis must demonstrate that the proposed score changes the information structure, interpretation or empirical utility rather than simply reversing an existing scale.

\section{Intersectionality and the Limits of Additive Social Models}

The literature on social inequality further complicates the assumption that social determinants can simply be added together. Intersectionality draws attention to the ways in which social positions and systems of power interact, producing experiences that may not be adequately represented by examining each dimension separately. \cite{crenshaw_1991,bowleg_2008}

The implication for measurement is not that every possible interaction must be included in the primary latent model. Such a strategy would rapidly become statistically unstable and conceptually opaque. The stronger implication is that the meaning of the proposed indicators may vary across groups. Income, education, housing, race, sex, social support and access to care do not necessarily operate identically in every population.

Accordingly, equity is a measurement question as well as a prediction question. A score can predict equally well across groups while measuring different latent quantities, and a score can demonstrate scalar differences while exhibiting non-invariance in the relationship between indicators and the latent construct. The thesis therefore needs measurement-invariance analysis before interpreting group comparisons as substantive differences in health-related potential.

There is also an ethical distinction between measuring constraint and assigning blame. If inadequate transport or unstable housing lowers an individual's H-potential score, the interpretation should be that the person's health-related capability is constrained by context, not that the individual possesses an intrinsic personal deficit. This distinction becomes essential if the metric is ever translated into clinical prioritisation, resource allocation or policy targeting.

\section{The Measurement Model Is Part of the Theory}

The preceding literatures converge on a methodological point that is often treated as secondary: the statistical form of the measurement model expresses a claim about what the construct is.

A reflective model assumes that an underlying latent variable generates covariance among its indicators. Removing one indicator should therefore not fundamentally redefine the construct, although precision may change. A formative or composite model instead treats indicators as contributors to the construct. Indicators may therefore be weakly correlated or even conceptually heterogeneous while remaining jointly necessary. The distinction is especially consequential when biological capacities are combined with social conditions.

The intrinsic-capacity literature itself has not resolved this issue. The scoping review of the field found that most studies did not clearly define the measurement model and that both reflective and formative interpretations appear in the literature. \cite{who_intrinsic_capacity_2017,beard_intrinsic_2019} This makes it inappropriate for the present thesis to state that SEM has validated H-potential simply because fit indices are acceptable.

The empirical programme must therefore distinguish at least three questions. First, is there coherent covariance within each domain? Second, are the domains empirically related strongly enough to justify a higher-order construct? Third, does collapsing them into a scalar score retain rather than destroy substantive information? These questions lead naturally to competing models rather than a single confirmatory endpoint.

The planned comparisons are therefore: (i) a correlated-domain model; (ii) a second-order model in which a general factor explains covariance among domains; and, where identification and sample size permit, (iii) a bifactor or related general-specific representation. Model choice should consider substantive interpretation, parameter stability, local fit, factor determinacy, discriminant validity and measurement invariance. Information criteria may assist comparison, but the statistically best-fitting model is not automatically the theoretically correct model.

The same logic applies to indicator selection. Modification indices should not become a mechanism for discovering the theory after seeing the data. Correlated residuals, cross-loadings and item removals require substantive justification and should be treated as sensitivity analyses where possible. Otherwise the model risks becoming a description of idiosyncrasies in one sample rather than a measurement theory capable of replication.

Misspecifying formative indicators as reflective, or the reverse, can bias
structural relations and produce apparently satisfactory models with the wrong
causal interpretation. \cite{jarvis_2003} Composite construction adds further
normative choices: normalisation, weighting and aggregation determine which
disadvantages can compensate for others and can materially alter rankings.
Equal weights are not the absence of judgement; they are one particular
judgement. \cite{greco_2019}

For that reason, fit cannot choose the construct on its own. Content validity
determines what must be represented; measurement theory determines which arrows
are defensible; sensitivity analysis shows whether conclusions depend on
weights or substitution assumptions; and external validation tests whether the
representation travels. A single high-fitting SEM followed by a factor score
would reproduce, rather than resolve, the problem.

\section{The Unresolved Space and the Specific Contribution of H-Potential}

The literature does not reveal one missing variable. It reveals a set of
interfaces that existing traditions handle separately. Clinical measures are
often strongest when describing disease, impairment or realised function.
Intrinsic capacity shifts attention towards positive capacity but remains
contested in its operationalisation. Capability theory explains why opportunity
cannot be inferred from functioning alone, yet complete opportunity sets are
difficult to observe. Social-determinants research explains how opportunities
are structured, but it does not by itself yield an individual health
measurement model. Frailty and social-vulnerability indices predict risk well,
but generally organise biological and social disadvantage as accumulated
vulnerability. QALYs and DALYs support comparison and valuation, but do not
model the conversion of capacity into functioning.

H-Potential is proposed at these interfaces. Its first gap is
\emph{conceptual}: it distinguishes available biological capacity, enabling
conditions, structural constraints, realised function and downstream outcomes
within one ordered account. Its second gap is \emph{measurement-related}: it
does not assume that all three domains are reflective manifestations of one
cause, and it permits a profile or correlated-domain result when a total score
is indefensible. Its third gap is \emph{functional-form related}: additive,
interactive, multiplicative and non-compensatory aggregations are treated as
competing hypotheses rather than hidden design choices. Its fourth gap is
\emph{evaluative}: the construct is tested against intrinsic capacity, frailty,
social vulnerability and current health before any claim of added value. Its
fifth gap is \emph{interpretive}: structural constraint remains visible as
context rather than being redescribed as an intrinsic personal deficit.

These are proposed contributions, not established achievements. Each has a
failure condition. The conceptual ordering fails if the domains cannot be
distinguished. The measurement contribution fails if the indicators lack
content validity or the structure cannot be replicated. The aggregation
contribution fails if rankings depend arbitrarily on scaling or weights. The
incremental contribution fails if established measures contain the same
information more parsimoniously. The interpretive contribution fails if users
reduce the output to a decontextualised rank. By stating these conditions, the
theory becomes elaborated enough to be criticised and precise enough to be
tested.

The resulting novelty claim is therefore intentionally conditional. The thesis
does not claim discovery of multidimensional health, social determination,
capability, latent measurement or prognosis. It asks whether their theoretically
specified interface can be represented without collapsing capacity into
function, context into pathology, or prediction into construct validity. That
is narrower than a universal theory of health, but it is more defensible and
more useful as a doctoral contribution.

\section{What the Literature Supports---and What It Does Not}

The literature supports several propositions strongly enough to motivate empirical testing. First, health in later life is multidimensional. Second, health state, functional ability, intrinsic capacity, capability and vulnerability are related but non-identical concepts. Third, intrinsic capacity has become an important organising framework for healthy ageing, but its measurement remains heterogeneous. Fourth, social and environmental conditions influence the extent to which biological capacity can be translated into functioning. Fifth, social vulnerability and frailty already capture substantial portions of the biological and contextual risk space. Sixth, multidimensional health metrics can be constructed using latent-variable and item-response methods, but their validity depends on explicit conceptual and measurement assumptions.

The literature does not support several stronger claims. It does not support the assertion that multidimensional health measurement is absent. It does not support treating the proposed algebraic expression as an identified empirical estimator. It does not support assuming that social conditions are interchangeable reflective indicators of a latent biomedical state. It does not support assuming that a general factor exists because the theory contains several domains. Finally, prognostic association does not establish causal leverage: a variable can predict mortality without being an intervention target, and a social burden can be predictive without being a component of a reflective latent trait.

These limitations are not obstacles to the thesis; they define its scientific task. The research question is therefore whether the proposed construct occupies a distinguishable measurement space between intrinsic capacity, frailty, social vulnerability and conventional health status, and whether that distinction carries information about subsequent outcomes.

\section{The H-potential Equation as a Heuristic, Not an Estimator}

The original formulation
\begin{equation}
H_{\text{potential}}=(L\cdot Q)S-B
\end{equation}
was useful because it expressed an intuitive proposition: biological potential is modified by enabling conditions and constrained by structural burden. The literature, however, does not justify the equation as an empirical law.

Multiplication by $S$ assumes that social conditions proportionally scale every unit of biomedical potential. Subtraction of $B$ assumes that systemic burden exists on a common additive scale with the quantity produced by $(L\cdot Q)S$. Neither assumption follows from WHO's framework, capability theory, social-determinants research or frailty science. In particular, capability theory does not imply that a reduction in social opportunity should be numerically equivalent to a unit reduction in biological capacity.

The equation should therefore remain as a conceptual heuristic while the empirical model is estimated independently. The thesis should not interpret a standardised SEM factor score as years of healthy life, nor should it convert an arbitrary latent scale into a clinically meaningful unit without external calibration.

\section{Candidate Mathematical Representations and Their Tests}

Equation~\ref{eq:hp_general_class} defines a class of admissible models rather
than selecting one. Four candidates expose the principal aggregation assumptions:
\begin{align}
H_{P,i}^{(A)}&=\beta_0+\beta_1BP_i+\beta_2SM_i-\beta_3SB_i+\varepsilon_i,
\label{eq:hp_additive}\\
H_{P,i}^{(I)}&=\beta_0+\beta_1BP_i+\beta_2SM_i-\beta_3SB_i
+\beta_4(BP_iSM_i)+\varepsilon_i,
\label{eq:hp_interaction}\\
H_{P,i}^{(M)}&=BP_i^{\alpha}SM_i^{\beta}(1-SB_i)^{\gamma},
\qquad \alpha,\beta,\gamma>0,
\label{eq:hp_multiplicative}\\
H_{P,i}^{(C)}&=BP_i\left[\delta SM_i^{\rho}
+(1-\delta)(1-SB_i)^{\rho}\right]^{1/\rho},
\quad 0<\delta<1.
\label{eq:hp_ces}
\end{align}
The bounded complement $1-SB_i$ makes zero a complete contextual bottleneck.
All components must be directionally harmonised and scaled before comparison.

The additive model assumes constant marginal contributions and complete
compensation on a common scale. The interaction model tests whether enablement
changes the contribution of biological capacity; $\beta_4=0$ is the nested null
of no complementarity. The multiplicative model imposes positive monotonicity
and proportional trade-offs. The CES model varies substitutability:
$\rho\rightarrow1$ approaches a weighted arithmetic mean,
$\rho\rightarrow0$ the geometric case, and $\rho\rightarrow-\infty$ a
bottleneck limit. These are empirical and normative alternatives.

For the heuristic $(LQ)S-B$,
\begin{equation}
\frac{\partial H_P}{\partial L}=QS,
\qquad
\frac{\partial^2 H_P}{\partial L\,\partial S}=Q.
\label{eq:heuristic_derivatives}
\end{equation}
The positive cross-partial is a substantive complementarity claim: support
raises the marginal contribution of biological capacity. Additive barriers also
require commensurability between $B$ and $(LQ)S$. Failure of either assumption
rejects that functional form, not automatically the general construct.

The transition from construct to measurement is separate. For a reflective
domain $\eta_i$, indicator $j$ is represented as
\begin{equation}
x_{ij}=\nu_j+\lambda_j\eta_i+\epsilon_{ij},
\label{eq:reflective_measurement_bridge}
\end{equation}
whereas a formative composite reverses the direction,
$\eta_i=\sum_j w_jx_{ij}+\zeta_i$. The code therefore does not force social
and structural indicators into the reflective model. It compares admissible
measurement structures without outcomes and only then compares the candidate
realisations against prespecified comparators and downstream outcomes.

The more defensible conceptual representation is therefore
\begin{equation}
H_{\text{potential}}=f(BP,SM,SB),
\end{equation}
where $BP$ denotes Biomedical Potential, $SM$ Social Modifiers and $SB$ Systemic Burdens, and $f$ is intentionally left unspecified at the theoretical stage. The empirical analysis is then responsible for determining whether these domains share a common higher-order structure and whether a scalar representation is warranted.

\section{From Literature to Testable Propositions}

The literature synthesis generates the following propositions.

\begin{enumerate}
    \item \textbf{Multidimensionality proposition.} The indicators selected for the proposed construct will represent more than one empirically distinguishable domain of health-related capability rather than a single undifferentiated health state.
    \item \textbf{Capacity proposition.} Biomedical Potential will show substantial convergence with established measures of intrinsic capacity and current functioning while retaining information not reducible to either.
    \item \textbf{Context proposition.} Social Modifiers and Systemic Burdens will demonstrate associations with health-related functioning that are not fully explained by biomedical capacity alone.
    \item \textbf{Distinctiveness proposition.} The proposed construct will demonstrate incremental information relative to pre-specified frailty and social-vulnerability comparators; otherwise its justification as a distinct metric will be limited.
    \item \textbf{General-factor proposition.} A higher-order H-potential factor will be retained only if the data demonstrate sufficient commonality, factor determinacy and substantive coherence. A correlated-domain model remains a valid endpoint if those conditions are not met.
    \item \textbf{Equity proposition.} The measurement structure will be sufficiently invariant across major demographic and socioeconomic groups to support meaningful comparisons. Where invariance fails, the thesis will report the failure rather than treating observed group differences as direct evidence of unequal health potential.
    \item \textbf{Prognostic proposition.} If a general H-potential score is empirically defensible, it will provide prognostic information for mortality and/or hospitalisation beyond established health and vulnerability measures. This proposition is secondary to construct validity and will not be used to select the measurement model.
\end{enumerate}

Each proposition has a corresponding empirical obligation. The first requires exploratory and confirmatory dimensionality analysis; the second requires convergent and discriminant validity testing; the third requires domain-level and sensitivity analyses; the fourth requires nested and incremental prognostic comparisons; the fifth requires competing measurement structures and factor-determinacy assessment; the sixth requires measurement-invariance analysis; and the seventh requires independently specified survival and hospitalisation models with internal validation.

This mapping also establishes a direct link to the computational implementation. The code should not merely calculate a final score. It must preserve the distinction between theory generation, measurement evaluation, score estimation and outcome prediction. Every model-selection decision must therefore be traceable to a stated theoretical proposition or methodological criterion.

\section{Implications for the Empirical Design and Codebase}

The literature review changes the empirical design in several concrete ways. First, indicator selection must be documented as a theory-to-variable mapping. Each indicator should have a stated construct, domain, expected direction and justification from the literature. Second, candidate measurement models must be fitted before outcome variables are introduced into model selection. Third, the analysis must include explicit tests of alternative structures rather than only the preferred model.

Fourth, the code must interrogate the possibility that different domains have different measurement ontologies. Biomedical indicators may plausibly be reflective manifestations of underlying capacity, whereas social conditions may function as formative contributors, contextual variables or external causes. The analysis should therefore avoid forcing all indicators into one reflective CFA merely because a common SEM framework is available.

Fifth, the analysis must distinguish model fit from substantive validity. A model can fit well and still measure the wrong construct. Conversely, a theoretically coherent construct may show imperfect fit because the observed indicators are necessarily heterogeneous. Fit indices should therefore be interpreted alongside reliability, factor determinacy, construct validity, invariance and comparison with established instruments.

Sixth, the prognostic analysis must remain downstream of measurement. Mortality and hospitalisation should not be used to discover the latent structure. They should test whether the construct developed independently of the outcomes contains useful information. Comparisons with frailty, intrinsic capacity and social vulnerability are therefore not optional sensitivity analyses; they are central tests of the thesis's claim to added value.

Finally, the code should be designed to expose failure. It should be possible for the pipeline to conclude that no defensible general factor exists, that a domain is empirically redundant, that measurement invariance fails, or that prognostic performance adds little beyond established comparators. Such outcomes would not constitute programming failures. They would be scientifically informative findings about the proposed construct.

\section{Conclusion}

The literature reveals a mature field rather than an empty space. QALYs and DALYs address health loss and comparative valuation; health expectancy addresses expected population health over time; quality-of-life instruments describe experienced health; frailty captures accumulated vulnerability; intrinsic capacity describes the physical and mental capacities available to an individual; capability measures address opportunities to achieve valued states; and social-vulnerability measures quantify contextual disadvantage. The challenge is not to replace these traditions but to determine whether a coherent bridge among them is empirically defensible.

The proposed H-potential therefore enters the literature with a deliberately modest claim. Its novelty, if established, will not be that it is the first multidimensional health metric. Nor will it be that it includes social determinants or uses latent-variable methods. The potential contribution is the explicit testing of a health-bounded construct in which biological capacity, enabling social conditions and structural constraints are considered jointly, while measurement validity is kept analytically separate from prognostic performance.

That proposition remains open. If the data show that H-potential is indistinguishable from intrinsic capacity, frailty, social vulnerability or an established health metric, the thesis must conclude that the proposed integration does not warrant a separate construct. If the domains are coherent but a general factor is not, the appropriate result may be a multidomain profile rather than a scalar score. If the score is coherent and incrementally predictive, the contribution becomes stronger. In all cases, the scientific value lies in allowing the empirical evidence to determine which interpretation survives interrogation.
